\documentclass[../../../main.tex]{subfiles}

\begin{document}

\section{Advanced Inequalities}




Muirhead's inequality is one of the powerful proving tools that is
also a generalization of AM-GM inequality.

\begin{theorem}{Muirhead's Inequality}{Muirhead's Inequality}
Let $(a_i)$ and $(b_i)$ be sequences such that $(a_i) \succ (b_i)$.
Then, the following inequality holds for all nonnegative $x_i$.
\[
\sum_\text{sym} {x_1}^{a_1} \cdots {x_n}^{a_n} 
\geq \sum_\text{sym} {x_1}^{b_1} \cdots {x_n}^{b_n}
\]
The equality holds if and only if $(a_i)$ and $(b_i)$ are identical
or if all $x_i$ are identical, which is the most common case.
\end{theorem}






Another very powerful inequality is Shur's inequality.
Shur's inequality comes with many alternative forms, and it is
generally encouraged to prove the derivation as there are also
not-so-famous ones.

\begin{theorem}{Schur's Inequality}{Schur's Inequality}
For $a, b, c \geq 0$ and $r \geq 0$,
\begin{align*}
    \sum_\text{cyclic} a^r(a - b)(a - c)
    &\geq 0 \\
    a^r(a - b)(a - c) + b^r(b - a)(b - c) + c^r(c - a)(c - b)
    &\geq 0
\end{align*}
holds with equality if $a = b = c$.
\end{theorem}
\begin{proof}
First, notice that the left hand side of the inequality is cyclic.
Therefore without loss of generality, let $a \geq b \geq c \geq 0$.
Rewriting the inequality,
\begin{align*}
    a^r(a - b)(a - c) + b^r(b - a)(b - c) + c^r(c - a)(c - b)
    &\geq 0 \\
    (a - b)\left[ a^2(a - c) - b^r(b - c) \right] + c^r(c - a)(c - b)
    &\geq 0
\end{align*}
holds. Moreover by construction, $a - b$, $a^2(a - c) - b^r(b - c)$,
and $c^r(c - a)(c - b)$ are all greater than or equal to zero.
Since the sum of all non-negative numbers is a non-negative number,
the inequality holds.
\end{proof}

It won't be an exaggeration to say that the real power of Schur's
inequality in olympiads come from its corollaries.

\begin{corollary}{}{Shur's Inequality Corollary I}
The following inequality holds for all reals $a, b, c$.
\[
a^2 + b^2 + c^2 \geq ab + bc + ca
\]
\end{corollary}
\begin{proof}
Substituting $r = 0$ in Schur's inequality, the following is obtained.
\begin{align*}
    (a - b)(a - c) + (b - a)&(b - c) + (c - a)(c - b)
    \geq 0 \\
    \therefore a^2 + b^2 + c^2
    &\geq ab + bc + ca
\end{align*}
However, notice that unlike the condition for $a, b, c$ in Schur's
inequality, the inequality above holds for negative $a, b, c$.
Without loss of generality, let $a \geq b \geq c$. If $c < 0$ and
$a, b \geq 0$, then it is self evident that the inequality is holds.
Similarly, if two variables are negative, the inequality also holds.
Lastly, when all $a, b, c$ are all negative, then the inequality
holds as positive numbers are greater than negative sum.
\end{proof}






Minkowski's inequality is a generalization of the triangle inequality.

\begin{theorem}{Minkowski's Inequality}{Minkowski's Inequality}
For $a_i, b_i > 0$ and $p > 1$, the following inequality is holds.
\[
\left( \sum_{k = 1}^n (a_k + b_k)^p \right)^\frac{1}{p}
\leq
\left( \sum_{k = 1}^n a_k^p \right)^\frac{1}{p}
    + \left( \sum_{k = 1}^n b_k^p \right)^\frac{1}{p}
\]
The equality holds if and only if there exists $\lambda > 0$ such that
$a_i = \lambda b_i$ for all $1 \leq i \leq n$.
\end{theorem}


* Bernoulli’s Inequality
* Equal Value Principle

\subsection{Symmetric and Cyclic Inequalities}
* Newton's Inequality
* Maclaurin's Inequality
* Popoviciu's Inequality

Nesbitt's inequality can be considered as the direct consequence
of Cauchy inequality.

\begin{theorem}{Nesbitt's Inequality}{Nesbitt's Inequality}
For arbitrary $a, b, c > 0$, the following inequality holds.
\[
\sum_\text{cyc} \frac{a}{b + c} \geq \frac{3}{2}
\]
\end{theorem}
\begin{proof}
Consider the following inequality by Cauchy inequality.
\begin{align*}
    \left( (a + b) + (b + c) + (c + a) \right) &\left(
        \frac{1}{a + b} + \frac{1}{b + c} + \frac{1}{c + a}
    \right) \\
    &\qquad\qquad\qquad\quad
    \geq (1 + 1 + 1)^2
    = 9
\end{align*}
Therefore, the following inequalities hold.
\begin{align*}
    &\quad\,\ 2 \left(
        \frac{a + b + c}{a + b} + \frac{a + b + c}{b + c}
            + \frac{a + b + c}{c + a}
    \right) \\
    &= 2 \left(
        \frac{c}{a + b} + \frac{a}{b + c} + \frac{b}{c + a} + 3
    \right)
    \geq 9
\end{align*}
Thus, the inequality holds.
\end{proof}

Nesbitt's inequality is less frequently used than other inequalities,
but when applied, it can be a very powerful inequality.



\end{document}

